\documentclass[11pt]{article}
\usepackage{amsmath,amssymb,mathtools}
\usepackage{fontspec,unicode-math}
\setmainfont{Latin Modern Roman}
\setmathfont{Latin Modern Math}
\usepackage{xurl,hyperref}
\newcommand{\sref}[2]{\hyperlink{source:#1:#2}{[S#2]}}
\newcommand{\cataloguescope}[1]{\paragraph{Mathematical statement.}#1}
\begin{document}
\section{Research Problems in Function Theory — Problem 7.35}\label{statement}

\subsection{Definitions and mathematical statement}

\paragraph{Shared notation.}$\mathbb N=\{1,2,\ldots\}$, and $\mathbb Z,\mathbb Q,\mathbb R,\mathbb C$ denote the integers, rationals, reals and complex numbers. Any other conventions are specified locally.


On the unit circle use angular Lebesgue measure. A real integrable $u$ is in $\operatorname{BMO}$ if
\[\sup_I\frac1{|I|}\int_I|u-u_I|\,d\theta<\infty,\qquad
u_I=|I|^{-1}\int_Iu\,d\theta,
\]
where $I$ ranges over arcs. For a real bounded function $b$, its conjugate $\widetilde b$ is the circular Hilbert transform with zero mean: its Fourier coefficient of index $k\ne0$ is $-i\operatorname{sgn}(k)\widehat b(k)$ and its coefficient of index zero is zero. Equalities of boundary functions are taken almost everywhere and $\|\cdot\|_\infty$ is the essential supremum.

\paragraph{Statement recorded in the supplied source.}
Given real $u\in\operatorname{BMO}$, determine the optimal value
\[\inf\{\|b_2\|_\infty:u=b_1+\widetilde b_2,\quad b_1,b_2\text{ real and bounded}\}.\] \sref{statement}{2}
Fefferman's decomposition assures admissible pairs exist. The original relates the optimisation to Baernstein's zero-free univalent-function factorisation question in Problem5.58. \sref{statement}{2}
\subsection{Short English statement}
Given a real function of bounded mean oscillation, determine the smallest possible size of its bounded conjugate component in a bounded-plus-conjugate decomposition.
\subsection{Sources}
\begin{description}
\item[\hypertarget{source:statement:1}{S1}] ULAM.AI and the UnsolvedMath Contributors, \emph{UnsolvedMath: A Curated Collection of Open Mathematics Problems}, catalogue entry. CC BY4.0. \url{https://huggingface.co/datasets/ulamai/UnsolvedMath}.
\item[\hypertarget{source:statement:2}{S2}] W. K. Hayman and E. F. Lingham, \emph{Research Problems in Function Theory} (2018), Chapter7, Problem7.35 and its complete update; Problem5.58 and its update. \url{https://arxiv.org/abs/1809.07200}.
\end{description}
\label{end:statement}
\end{document}
